%! Logarithmic Function Context and Data Modeling — Unit 5, Lesson 5.7 — Slides (Beamer)
\documentclass[aspectratio=169, 11pt]{beamer}
\usepackage{precalculus-beamer}   % provides \navyheader{} and \sectionlabel[color]{}

\begin{document}

% TITLE SLIDE
{
\setbeamercolor{background canvas}{bg=navy}
\begin{frame}[plain]
  \vspace{2.8cm}\hspace{0.55cm}%
  \begin{minipage}{0.9\textwidth}
    {\color{goldacc}\bfseries\small LESSON 2.14}\\[0.5cm]
    {\color{white}\bfseries\LARGE Logarithmic Function Context and Data Modeling}\\[1.2cm]
    {\color{skymid}\small Precalculus~$\cdot$~Unit 5: Logarithmic Functions}
  \end{minipage}
\end{frame}
}

% SLIDE 1 — HOOK
\begin{frame}[t, plain]
  \navyheader{Hook: When Does Growth Slow Down?}

  \vspace{0.3cm}
  \textbf{A scientist observes:}
  \begin{itemize}
    \setlength{\itemsep}{8pt}
    \item Bacteria double in the \textbf{1st hour}.
    \item They double again after \textbf{2 more hours}.
    \item They double again after \textbf{4 more hours}.
  \end{itemize}

  \vspace{0.4cm}
  \begin{block}{Think About It}
    As the population grows, the time needed for each doubling \textbf{increases}.\\[6pt]
    If we let the output = number of doublings and input = elapsed time, what kind
    of function has \emph{equal output increments} for \emph{proportionally growing inputs}?
  \end{block}
\end{frame}

% SLIDE 2 — WHEN TO USE A LOG MODEL
\begin{frame}[t, plain]
  \navyheader{When to Use a Logarithmic Model (EK 2.14.A.1)}

  \vspace{0.3cm}
  \begin{block}{Key Idea}
    A \textbf{logarithmic model} is appropriate when input values change
    \textbf{proportionally} over equal-length output intervals.
  \end{block}

  \vspace{0.3cm}
  \renewcommand{\arraystretch}{1.6}
  \begin{tabular}{@{} l @{\quad} l @{}}
    \textbf{Exponential:} & outputs multiply by a constant over equal \emph{input} intervals \\
    \textbf{Logarithmic:} & inputs multiply by a constant over equal \emph{output} intervals \\
  \end{tabular}

  \vspace{0.4cm}
  \textbf{Example table:}

  \vspace{0.2cm}
  \begin{center}
  \renewcommand{\arraystretch}{1.5}
  \begin{tabular}{|c|c|c|c|c|}
  \hline
  $x$ & $2$ & $4$ & $8$ & $16$ \\
  \hline
  $y$ & $1$ & $2$ & $3$ & $4$ \\
  \hline
  \end{tabular}
  \end{center}

  \vspace{0.2cm}
  As $y$ increases by 1, $x$ \textbf{doubles} $\;\Rightarrow\;$ proportion = 2
  $\;\Rightarrow\;$ model: $y = \log_2(x)$
\end{frame}

% SLIDE 3 — BUILDING A MODEL FROM TWO PAIRS
\begin{frame}[t, plain]
  \navyheader{Building a Model from Two Input-Output Pairs (EK 2.14.A.2)}

  \vspace{0.3cm}
  \textbf{General form:} $f(x) = a\log_b x + k$

  \vspace{0.3cm}
  \textbf{Worked example:} Given $f(1) = 0$ and $f(e) = 3$.

  \begin{enumerate}
    \setlength{\itemsep}{6pt}
    \item Substitute $(1,\,0)$: $\;0 = a\ln(1) + k = 0 + k \;\Rightarrow\; k = 0$
    \item Substitute $(e,\,3)$: $\;3 = a\ln(e) = a \cdot 1 \;\Rightarrow\; a = 3$
    \item \textbf{Model:} $f(x) = 3\ln(x)$
  \end{enumerate}

  \vspace{0.3cm}
  \begin{block}{System of Equations Strategy}
    \begin{enumerate}
      \setlength{\itemsep}{4pt}
      \item Plug both pairs into $f(x) = a\log_b x + k$
      \item Use $\log_b(1) = 0$ to find $k$ from the first equation
      \item Substitute $k$ into the second equation and solve for $a$
    \end{enumerate}
  \end{block}
\end{frame}

% SLIDE 4 — TRANSFORMATIONS
\begin{frame}[t, plain]
  \navyheader{Transformations of a Logarithmic Model (EK 2.14.A.3)}

  \vspace{0.3cm}
  \textbf{General form:} $f(x) = a\log_b(x - h) + k$

  \vspace{0.2cm}
  \begin{itemize}
    \setlength{\itemsep}{6pt}
    \item $a$ --- vertical dilation (stretch/compression, and reflection if $a < 0$)
    \item $h$ --- horizontal shift ($h > 0$: right; $h < 0$: left)
    \item $k$ --- vertical shift
  \end{itemize}

  \vspace{0.3cm}
  \begin{block}{Decibel Example}
    Sound level: output = 0 dB at $x = 1$; increases 10 dB per tenfold increase in $x$.\\[6pt]
    $\Rightarrow\; k = 0$, base $b = 10$, dilation $a = 10$\\[4pt]
    $S(x) = 10\log_{10}(x)$
  \end{block}

  \vspace{0.2cm}
  The full decibel formula $S(I) = 10\log_{10}(I/I_0)$ is a \textbf{vertical shift}
  of this model by $+120$ (since $I_0 = 10^{-12}$).
\end{frame}

% SLIDE 5 — NATURAL LOG MODELS AND PREDICTIONS
\begin{frame}[t, plain]
  \navyheader{Natural Logarithm Models and Predictions (EK 2.14.A.5, 2.14.A.6)}

  \vspace{0.3cm}
  \begin{block}{Natural Log}
    $\ln x = \log_e x$ \qquad Key identities: $\ln(e^n) = n$ \quad and \quad $e^{\ln x} = x$
  \end{block}

  \vspace{0.3cm}
  \textbf{Evaluating:} $f(x) = 3\ln x$

  \vspace{0.1cm}
  \begin{itemize}
    \setlength{\itemsep}{6pt}
    \item $f(e^4) = 3\ln(e^4) = 3 \cdot 4 = \mathbf{12}$
    \item $f(1) = 3\ln(1) = 3 \cdot 0 = \mathbf{0}$
  \end{itemize}

  \vspace{0.3cm}
  \textbf{Solving (inverting):} Find $x$ when $f(x) = 9$
  \[
    3\ln x = 9 \;\Rightarrow\; \ln x = 3 \;\Rightarrow\; x = e^3
  \]

  \vspace{0.2cm}
  \textbf{Technology path (EK 2.14.A.4):} Logarithmic regression on a calculator fits
  $y = a + b\ln x$ to a data set automatically --- always interpret the result in context.
\end{frame}

% SLIDE 6 — CLASS PRACTICE
\begin{frame}[t, plain]
  \navyheader{Class Practice}

  \vspace{0.3cm}
  Work with your partner. Be ready to share.

  \vspace{0.3cm}
  \begin{enumerate}
    \setlength{\itemsep}{10pt}
    \item A data set: $x = 3, 9, 27, 81$ with $y = 1, 2, 3, 4$. Identify the proportion,
          write the model, and predict $y$ when $x = 243$.

    \item Given $g(1) = 5$ and $g(e^2) = 11$, find $a$ and $k$ for $g(x) = a\ln(x) + k$.
          Write the model and evaluate $g(e^4)$.

    \item A model is $h(x) = 2\log_{10}(x-1)+1$. State the domain. Find $h(101)$.
          Solve $h(x) = 5$ for $x$.
  \end{enumerate}
\end{frame}

\end{document}
